\chapter{Soluzione Proposta}
\label{ch:soluzione}

In questo capitolo viene descritta in dettaglio la soluzione realizzata, partendo dall'architettura complessiva del sistema per poi approfondire il funzionamento logico di ciascun componente e i principali dettagli implementativi. La piattaforma ChillChain si compone di due macro-funzionalità integrate: un marketplace per il trasporto condiviso refrigerato e un sistema di monitoraggio in tempo reale con rilevazione automatica delle anomalie.

% ============================================================================
\section{Architettura del sistema}
\label{sec:architettura}

L'architettura di ChillChain segue il paradigma a microservizi, in cui ciascun componente è responsabile di una funzione specifica e comunica con gli altri attraverso interfacce ben definite. La Figura~\ref{fig:architettura} illustra schematicamente i servizi e i flussi di comunicazione.

\begin{figure}[ht]
\centering
% Inserire qui il diagramma architetturale
% \includegraphics[width=\textwidth]{img/architettura.png}
\fbox{\parbox{0.9\textwidth}{\centering\vspace{2cm}\textit{Diagramma architetturale del sistema}\vspace{2cm}}}
\caption{Architettura a microservizi di ChillChain. Le frecce continue rappresentano comunicazioni HTTP/REST, le frecce tratteggiate i flussi MQTT.}
\label{fig:architettura}
\end{figure}

Il sistema è composto da sei servizi containerizzati, orchestrati tramite Docker Compose:

Il \textbf{backend Spring Boot} (porta 8081) rappresenta il nucleo applicativo: espone le API REST per la gestione di veicoli, tratte, spedizioni, utenti e notifiche, implementa l'autenticazione JWT e la logica di autorizzazione basata sui ruoli, sottoscrive i topic MQTT per ricevere telemetria e risultati dell'anomaly detection, e persiste tutti i dati su MongoDB.

Il \textbf{frontend React} (porta 5173) fornisce l'interfaccia utente, differenziata per ruolo: l'amministratore gestisce la flotta e le tratte, il cliente cerca e prenota spedizioni, il tecnico monitora i veicoli e riceve le notifiche di anomalia.

Il \textbf{broker MQTT Mosquitto} (porta 1883) gestisce la comunicazione asincrona tra i servizi Python e il backend, instradando i messaggi attraverso topic gerarchici.

Il \textbf{servizio di streaming} (porta 8002) simula i sensori a bordo dei veicoli, leggendo dataset CSV e pubblicando i dati telemetrici sul broker MQTT con una cadenza di un messaggio al secondo.

Il \textbf{servizio di anomaly detection} (porta 8003) sottoscrive la telemetria pubblicata dallo streamer, la analizza con il modello LSTM Autoencoder e pubblica i risultati su un topic dedicato.

Infine, \textbf{MongoDB} (porta 27017) funge da database centralizzato, organizzato in sei collezioni: \texttt{user}, \texttt{vehicle}, \texttt{trip}, \texttt{shipment}, \texttt{telemetry} e \texttt{notification}.

La comunicazione tra i servizi avviene su due canali distinti: le API REST, utilizzate per le interazioni sincrone tra frontend e backend, e il protocollo MQTT, impiegato per i flussi di dati telemetrici che richiedono comunicazione asincrona e disaccoppiamento tra produttori e consumatori. Questa separazione consente di avviare, fermare o sostituire i servizi Python senza impattare sul funzionamento del backend o del frontend.


% ============================================================================
\section{Gestione degli utenti e autenticazione}
\label{sec:autenticazione}

Il sistema prevede tre ruoli distinti, ciascuno con un perimetro di azioni e una vista dell'interfaccia dedicati:

\begin{itemize}
    \item \textbf{ADMIN}: gestisce la flotta (aggiunta e rimozione di veicoli), visualizza e cancella tratte e spedizioni, avvia e interrompe le simulazioni di telemetria.
    \item \textbf{CLIENT}: ricerca tratte compatibili con le proprie esigenze di spedizione, confronta le opzioni disponibili (veicoli vuoti e tratte condivise) e prenota la spedizione.
    \item \textbf{TECHNICIAN}: accede alla dashboard di monitoraggio, visualizza lo storico telemetrico dei veicoli e gestisce le notifiche di anomalia.
\end{itemize}

L'autenticazione è implementata tramite token JWT (\textit{JSON Web Token}) con sessioni stateless. All'atto del login, il backend verifica le credenziali --- la password è memorizzata con hashing BCrypt --- e genera un token firmato contenente l'email dell'utente e il suo ruolo. Il token viene incluso nell'header \texttt{Authorization} di ogni richiesta successiva e validato da un filtro che precede l'elaborazione della richiesta stessa. La catena di sicurezza Spring Security definisce quali endpoint sono accessibili a ciascun ruolo: ad esempio, le operazioni di gestione della flotta richiedono il ruolo \texttt{ADMIN}, mentre l'accesso alla telemetria e alle notifiche è riservato al ruolo \texttt{TECHNICIAN}. Gli endpoint di registrazione, autenticazione e ricerca tratte sono invece pubblici, così da consentire l'onboarding autonomo dei clienti.


% ============================================================================
\section{Marketplace: ricerca e prenotazione delle tratte}
\label{sec:marketplace}

Il cuore commerciale della piattaforma è il meccanismo che consente ai clienti di cercare tratte compatibili con le proprie esigenze e di prenotare spedizioni, eventualmente condividendo il veicolo con altri clienti. Il flusso si articola in cinque fasi.

\subsection{Fase 1 --- Ricerca dei veicoli disponibili}

Quando un cliente inserisce i dati della propria spedizione (indirizzo di partenza e arrivo, dimensioni, peso, necessità di refrigerazione e data di arrivo desiderata), il backend esegue una prima query per individuare i veicoli della flotta che soddisfano i requisiti dimensionali e non hanno ancora alcuna tratta assegnata. Questa query è implementata come aggregation pipeline su MongoDB: il primo stadio filtra i veicoli per peso massimo, dimensioni e tipologia di refrigerazione; il secondo esegue un \texttt{\$lookup} sulla collezione \texttt{trip}; il terzo seleziona solo i veicoli il cui array di tratte risulta vuoto. Per ciascun veicolo disponibile, il sistema invoca la Routes API di Google Maps per calcolare il percorso tra la partenza e l'arrivo indicati dal cliente, ottenendo distanza, durata e polyline codificata. Il prezzo viene calcolato come $\text{prezzo} = \text{distanceKm} \times \text{pricePerKm}$, dove \texttt{pricePerKm} è il costo chilometrico associato al veicolo.

\subsection{Fase 2 --- Ricerca delle tratte condivise}

Parallelamente, il sistema cerca tratte già programmate che abbiano capacità residua sufficiente e data di arrivo compatibile. Anche questa ricerca è implementata come aggregation pipeline: il primo stadio esegue un \texttt{\$lookup} con la collezione \texttt{vehicle} per recuperare la tipologia di refrigerazione; il secondo filtra per capacità residua, refrigerazione e data; il terzo esegue un ulteriore \texttt{\$lookup} con la collezione \texttt{shipment} per calcolare il prezzo scontato come $\text{price} = \min(\text{shipments.price}) \div 2$.

\subsection{Fase 3 --- Validazione geografica delle tratte condivise}

Le tratte condivise individuate nella fase precedente vengono ulteriormente filtrate verificando la compatibilità geografica con la spedizione richiesta. Per ciascuna tratta candidata, il sistema decodifica la polyline del percorso attuale e calcola la distanza geodetica (formula di Haversine) tra i punti di ritiro e consegna del nuovo pacco e il tracciato della tratta. Se entrambi i punti distano meno di 1,5~km dalla polyline, la tratta è considerata geograficamente compatibile. Un ulteriore controllo verifica che il punto di ritiro preceda quello di consegna lungo il percorso, impedendo deviazioni illogiche.

\subsection{Fase 4 --- Ricalcolo del percorso con waypoint}

Per le tratte che superano la validazione geografica, il sistema ricalcola il percorso includendo i waypoint intermedi di tutti gli shipment (sia quelli già assegnati sia quello nuovo). L'ordinamento dei waypoint è affidato a un algoritmo greedy con vincoli: partendo dalla posizione corrente, seleziona il waypoint disponibile più vicino, rispettando il vincolo che il punto di ritiro di ciascuna spedizione debba essere visitato prima del corrispondente punto di consegna. Un delivery diventa ``disponibile'' solo dopo che il relativo pickup è stato inserito nella sequenza. Il percorso ricalcolato viene confrontato con quello originale: se la differenza di durata non supera i 30~minuti, la tratta viene inclusa tra i risultati.

\subsection{Fase 5 --- Pricing con divisore esponenziale}

Il sistema di pricing incentiva la condivisione attraverso un divisore esponenziale. Se il veicolo ospita già $n$ spedizioni, il prezzo per il nuovo cliente viene calcolato come:

\begin{equation}
\label{eq:pricing}
\text{prezzo} = \frac{\text{distanceKm}_{\text{shipment}} \times \text{pricePerKm}_{\text{vehicle}}}{2^n}
\end{equation}

Dove $\text{distanceKm}_{\text{shipment}}$ è la distanza del tratto specifico del nuovo cliente e $n$ è il numero di spedizioni già presenti. In questo modo, la prima spedizione paga il prezzo pieno, la seconda paga metà, la terza un quarto e così via. Questo meccanismo crea un incentivo economico crescente: più il veicolo è occupato, più il trasporto diventa conveniente per i nuovi clienti, incentivando il riempimento del mezzo e riducendo i costi per tutti i partecipanti.

\subsection{Selezione e conferma della tratta}

Il cliente visualizza tutte le opzioni --- veicoli vuoti e tratte condivise --- ordinate per prezzo, e può esaminare il percorso su mappa prima di confermare. Al momento della selezione, il backend esegue un'operazione transazionale: se la tratta è nuova (veicolo senza trip assegnato), viene creata una nuova entità \texttt{Trip} con la capacità residua decrementata delle dimensioni del pacco; se la tratta è condivisa, viene aggiornata l'entità \texttt{Trip} esistente con il nuovo percorso (inclusivo dei waypoint) e la capacità residua aggiornata. In entrambi i casi, viene creata un'entità \texttt{Shipment} associata alla tratta.


% ============================================================================
\section{Generazione dei dati telemetrici}
\label{sec:generazione_dati}

La simulazione della telemetria di bordo si basa su un modello termodinamico che riproduce il comportamento fisico di un impianto frigorifero a compressione di vapore. Il generatore, derivato dal dataset Kaggle \textit{Simulated Refrigerator Fault Diagnosis} \cite{kaggle_fridge} ed esteso per le esigenze del progetto, modella i seguenti fenomeni: lo scambio termico tra l'ambiente esterno e la cella frigorifera attraverso le pareti (parametrizzato dal coefficiente $UA_{\text{cab}}$), il ciclo di compressione con controllo proporzionale della velocità del compressore rispetto al setpoint di temperatura, il calcolo del COP (\textit{Coefficient of Performance}) basato sul ciclo di Carnot con fattore di efficienza, l'accumulo e lo scioglimento del ghiaccio sull'evaporatore, le aperture casuali della porta con carichi termici variabili e i cicli periodici di sbrinamento.

Il generatore produce dataset CSV contenenti 120 campioni (120 minuti con campionamento a 60 secondi), strutturati in due fasi: i primi 60 minuti in condizioni normali, seguiti da 60 minuti con un guasto iniettato. Sono previste 13 tipologie di guasto, tra cui: fouling del condensatore (lieve e severo), degradazione e guasto della ventola dell'evaporatore, sottocariamento e sovracariamento del refrigerante, deriva del sensore di temperatura, inefficienza del compressore, presenza di gas non condensabili e guasti compositi. Ciascun guasto è parametrizzato con un range stocastico, generando variabilità tra le simulazioni.

Il servizio di streaming legge questi CSV e pubblica ciascuna riga sul topic MQTT \texttt{fridge/\{vehicleName\}/telemetry} con un intervallo configurabile (default: 1 secondo), includendo un timestamp ISO e un indice di riga progressivo. Al termine del dataset, pubblica un messaggio di completamento con \texttt{stream\_status: "completed"}, che segnala agli altri servizi la fine dello stream.


% ============================================================================
\section{Pipeline di anomaly detection}
\label{sec:anomaly_detection}

La rilevazione delle anomalie si articola in tre fasi: l'addestramento offline del modello, l'inferenza in tempo reale e il meccanismo di smoothing per la riduzione dei falsi positivi.

\subsection{Addestramento del modello}

Il modello LSTM Autoencoder è stato addestrato su Google Colab con GPU Tesla T4, utilizzando un dataset composto esclusivamente da dati di funzionamento normale. L'architettura del modello è la seguente: un encoder LSTM con 15 feature in ingresso e hidden size pari a 16, uno strato lineare di compressione che riduce la rappresentazione a uno spazio latente di dimensione 8, due strati lineari che inizializzano hidden state e cell state del decoder a partire dalla rappresentazione latente, un decoder LSTM con hidden size 16 alimentato da una sequenza di zeri (la ricostruzione dipende interamente dallo stato latente) e uno strato lineare di output che proietta la dimensione nascosta nelle 15 feature originali.

L'addestramento è stato condotto per 100 epoche con batch size di 64 e finestre temporali di 30 campioni. L'ottimizzatore Adam con learning rate iniziale $10^{-3}$ è stato affiancato da uno scheduler \texttt{ReduceLROnPlateau} che dimezza il learning rate dopo 5 epoche senza miglioramento della validation loss. Il training utilizza mixed precision (\texttt{GradScaler}) per accelerare il calcolo su GPU. Un meccanismo di early stopping con pazienza di 10 epoche previene l'overfitting. Il dataset di training è stato suddiviso 80/20 in modo sequenziale (senza shuffle, per preservare la struttura temporale). La validation loss finale è risultata pari a circa 0,00398, con un tempo di addestramento complessivo di circa 715 secondi. Le 15 feature in ingresso vengono normalizzate mediante uno \texttt{StandardScaler} di scikit-learn, serializzato con Pickle e caricato a runtime dal servizio di inferenza.

La soglia di anomalia è stata determinata calcolando l'MSE di ricostruzione su tutto il dataset di training normale e prendendo il 95\textdegree{} percentile della distribuzione risultante ($\text{threshold} = 0,009868$). Questo approccio garantisce che il 95\% delle osservazioni normali produca un errore inferiore alla soglia, limitando il tasso di falsi positivi al 5\% teorico in condizioni nominali.

\subsection{Inferenza in tempo reale}

Il servizio di anomaly detection sottoscrive il topic \texttt{fridge/+/telemetry} tramite il broker MQTT, ricevendo i dati di tutti i veicoli. Per ciascun veicolo viene creato dinamicamente un'istanza di \texttt{VehicleAnomalyDetector} alla ricezione del primo messaggio, garantendo il supporto simultaneo di più veicoli con buffer e contatori indipendenti.

Ad ogni messaggio ricevuto, le 15 feature vengono estratte, normalizzate con lo scaler e inserite in un buffer a finestra scorrevole di 30 campioni. Quando il buffer è pieno, la finestra viene convertita in un tensore PyTorch e passata al modello in modalità di inferenza (\texttt{torch.no\_grad()}). L'errore di ricostruzione, calcolato come MSE tra input e output mediato su tutti i timestep e tutte le feature, viene confrontato con la soglia. Al completamento dello stream, il detector del veicolo viene rimosso e le sue risorse liberate.

\subsection{Smoothing con contatori consecutivi}

Per evitare che oscillazioni transitorie dell'errore di ricostruzione generino falsi allarmi, il sistema implementa un meccanismo di smoothing basato su contatori consecutivi. Vengono mantenuti due contatori per ciascun veicolo: un \texttt{anomaly\_counter}, che si incrementa ad ogni campione con $\text{MSE} > \text{threshold}$ e si azzera al primo campione normale, e un \texttt{normal\_counter}, che opera in modo simmetrico. La transizione dallo stato normale allo stato di anomalia avviene solo quando il contatore di anomalie supera la soglia di 20 campioni consecutivi; analogamente, il ritorno allo stato normale richiede 20 campioni consecutivi sotto soglia. Questo meccanismo introduce un ritardo controllato nella segnalazione (20 minuti con campionamento al minuto), ma elimina le notifiche spurie causate da picchi isolati, come quelli associati alle aperture della porta o ai cicli di sbrinamento.

Il risultato dell'analisi viene pubblicato sul topic \texttt{fridge/\{vehicleName\}/anomalies}, includendo l'errore di ricostruzione, lo stato dei contatori e, nelle sole transizioni di stato, un messaggio di allerta che il backend persisterà come notifica.


% ============================================================================
\section{Gestione delle notifiche}
\label{sec:notifiche}

Il backend sottoscrive il topic \texttt{fridge/+/anomalies} tramite un channel adapter Spring Integration dedicato. Alla ricezione di un messaggio di anomalia, il service activator verifica la presenza del campo \texttt{alertMessage}: se presente, indica una transizione di stato (da normale ad anomalia) e genera una nuova notifica persistente su MongoDB con i campi: nome del veicolo, timestamp, errore di ricostruzione, indice di riga e messaggio di allerta. Un meccanismo di fallback verifica se il messaggio di transizione è stato perso (ad esempio per un'interruzione temporanea della rete): se il flag \texttt{anomalyDetected} è attivo ma non esistono notifiche non lette per quel veicolo, ne viene creata una generica.

L'interfaccia del tecnico presenta le notifiche in ordine cronologico con un badge numerico che indica il conteggio delle notifiche non lette. Ciascuna notifica può essere contrassegnata come letta tramite un'azione dedicata, aggiornando lo stato sul database. Da una notifica, il tecnico può accedere direttamente alla dashboard telemetrica del veicolo interessato per analizzare lo storico dei dati e indagare la causa dell'anomalia.


% ============================================================================
\section{Dashboard telemetrica}
\label{sec:dashboard}

La dashboard di monitoraggio consente al tecnico di visualizzare lo storico telemetrico dei veicoli tramite grafici temporali interattivi. L'interfaccia è organizzata in due livelli: un selettore del veicolo, che elenca i veicoli con tratte attive, e un pannello di visualizzazione con grafici ad area realizzati tramite Recharts.

Per ciascun veicolo selezionato, il frontend richiede al backend l'elenco completo dei dati telemetrici memorizzati, ordinati cronologicamente. I dati vengono renderizzati in grafici con codifica cromatica basata sullo stato di guasto: i tratti in condizioni normali appaiono in verde, mentre quelli associati a valori anomali vengono evidenziati in rosso. Il tecnico può filtrare la visualizzazione per singolo parametro (ad esempio, visualizzare solo la temperatura della cella o solo la pressione di scarico) oppure visualizzare tutti i 15 parametri simultaneamente. L'asse temporale mostra il numero progressivo del campione, consentendo di individuare rapidamente il momento in cui il comportamento del sistema è cambiato.


% ============================================================================
\section{Deployment}
\label{sec:deployment}

L'intero sistema viene eseguito in container Docker orchestrati da Docker Compose. Il file di configurazione definisce i sei servizi, le porte esposte, le variabili d'ambiente e i volumi persistenti per MongoDB. La rete Docker interna consente ai servizi di comunicare tramite i nomi dei container: il backend raggiunge MongoDB tramite l'hostname \texttt{mongodb}, il broker MQTT tramite \texttt{mosquitto} e il servizio di streaming tramite \texttt{fridge-streamer}. L'avvio dell'intero sistema richiede un singolo comando:

\begin{verbatim}
docker-compose up --build
\end{verbatim}

Questa configurazione garantisce la riproducibilità dell'ambiente su qualsiasi macchina dotata del runtime Docker, eliminando le problematiche legate alla configurazione manuale delle dipendenze e delle connessioni di rete tra i servizi.